\begin{abstract}
In many computer vision and shape analysis tasks, practitioners are interested in learning from the \textit{shape} of the object in an image, while disregarding the object's orientation. To this end, it is valuable to define a \textit{rotation-invariant} representation of images, retaining all information about that image, but disregarding the way an object is rotated in the frame. To be practical for learning tasks, this representation must be \textit{computationally efficient} for large datasets and \textit{invertible}, so the representation can be visualized in image space. To this end, we present the selective disk bispectrum: a fast, rotation-invariant representation for image shape analysis. While the translational bispectrum has long been used as a translational invariant representation for 1D and 2D signals, its extension to 2D (disk) rotational invariance on images has been hindered by the absence of an invertible formulation and its cubic complexity. In this work, we derive an explicit inverse for the disk bispectrum, which allows us to define a ``selective" disk bispectrum, which only uses the minimal number of coefficients needed for faithful shape recovery. We show that this representation enables multi-reference alignment for rotated images—a task previously intractable for disk bispectrum methods. These results establish the disk bispectrum as a practical and theoretically grounded tool for learning on rotation-invariant shape data.
\end{abstract}
